\origpage{45}{}
\noindent by induction on $|V(G)| + |E(G)|$ that if $G$ is an intertwining of two $p$-linkages $L_1$ and $L_2$ then $G$ has rank $\le r(p)$.

\hspace*{12pt}(1) \textit{We may assume that $G$ has no isolated vertices.}

\textit{Subproof.} Suppose that $v \in V(G)$ has degree 0. For $i = 1, 2$, let $L'_i = L_i$ if $v \notin V(L_i)$, and $L'_i = L_i \setminus v$ if $v \in V(L_i)$. Then $G \setminus v$ is an intertwining of the $p$-linkages $L'_{\delm{i}\fix{1}}$ and $L'_2$\fixnote{fix:11.1-claim1}{Here $L'_i$ was just defined for $i = 1, 2$, and the two linkages are $L'_1$ and $L'_2$.}, and from our inductive hypothesis $G \setminus v$ has rank $\le r(p)$. But $G$ and $G \setminus v$ have the same rank, and the result follows. This proves (1).

\hspace*{12pt}(2) \textit{We may assume that $V(L_1) = V(L_2) = V(G)$.}

\textit{Subproof.} Suppose that $v \in V(G) \setminus V(L_2)$ say. By (1) there is an edge $e$ of $G$ incident with $v$, and since $L_1 \cup L_2 = G$ it follows that $e \in E(L_1)$. The graph $G/e$ (that is, obtained from $G$ by contracting $e$) is an intertwining of $L_1/e$ and $L_2$, if we regard $L_2$ as a subgraph of $G/e$ in the obvious way; and $L_1/e$ and $L_2$ are both $p$-linkages. \fix{We give the details. Let $e$ have ends $u, v$, and let $G/e$ be obtained from $G \setminus e$ by identifying $v$ with $u$. For a subgraph $K$ of $G$ with $v \notin V(K)$ or $e \in E(K)$, let $K/e$ be the corresponding subgraph of $G/e$. Then $L_2$, which contains neither $v$ nor $e$, is a linkage of $G/e$ with the same pattern, and $L_1/e$ is a linkage of $G/e$ whose pattern is obtained from that of $L_1$ by replacing $v$ by $u$, and which has no more terminals than $L_1$. We show that}

\fix{(*) \textit{Let $L \in \{L_1, L_2\}$, and let $M$ be a linkage of $G/e$ whose pattern is obtained from the pattern of $L$ by replacing $v$ by $u$. Then $M = K/e$ for some linkage $K$ of $G$ with the same pattern as $L$.}}

\fix{This implies that $G/e$ is an intertwining of $L_1/e$ and $L_2$, for if $M_1, M_2$ are linkages of $G/e$ with the patterns of $L_1/e, L_2$, then the corresponding $K_1, K_2$ have the patterns of $L_1, L_2$, so $K_1 \cup K_2 = G$, and $M_1 \cup M_2 = (K_1 \cup K_2)/e = G/e$. To prove (*), note first that if $u$ and $v$ are both terminals of $L$, then $\{u, v\}$ is a block of its pattern, since then $L = L_1$ and $e \in E(L_1)$. No edge of $M$ is incident in $G$ with both $u$ and $v$, since it would be a loop of $M$. Let $a$ and $b$ be the numbers of edges of $M$ incident with $u$ in $G/e$ that are incident in $G$ with $u$ and with $v$ respectively; so $u$ has degree $a + b \le 2$ in $M$, and $u$ is a terminal of $M$ if and only if $u$ or $v$ is a terminal of $L$. There are three ways to choose $K$:}

\fix{(A) if $b = 0$, let $K = M$;}

\fix{(B) if $a = 0$, let $K$ be $M$ with $u$ renamed $v$;}

\fix{(C) if $u \in V(M)$ and $a, b \le 1$, let $K$ be $M$ with $v$ and $e$ added.}

\fix{In each case $K$ is a linkage of $G$ with $K/e = M$ (in (C), $u$ and $v$ have degree $a + 1$ and $b + 1$ in $K$, and a circuit of $K$ would give one of $M$), and the pattern of $K$ becomes that of $M$ when $v$ is replaced by $u$. So the patterns of $K$ and $L$ agree after this replacement, and so they are equal if neither $K$ nor $L$ has $v$ as a terminal, or neither has $u$ as a terminal, or both have $\{u, v\}$ as a block, since in each of these cases the replacement is one-to-one on the blocks involved. We choose $K$ so that one of these holds.}

\fix{Suppose first that $b = 0$. If $v$ is not a terminal of $L$, we use (A), and then $v \notin V(K)$. If $u$ and $v$ are both terminals of $L$, then $\{u, v\}$ is a block of $L$, so $\{u\}$ is a block of $M$ and $a = 0$; we use (C), and then $\{u, v\}$ is a block of $K$. If $v$ is a terminal of $L$ and $u$ is not, then $u$ is a terminal of $M$, so $a \le 1$; we use (B) if $a = 0$ and (C) if $a = 1$, and in both cases $u$ is not a terminal of $K$. Now suppose that $b \ge 1$. If $a = 0$, and either $b = 2$ or $v$ is a terminal of $L$ and $u$ is not, we use (B); then $u \notin V(K)$, and $u$ is not a terminal of $L$ (if $b = 2$, because $u$ is not a terminal of $M$). Otherwise $a + b = 2$, or $a = 0$ and $b = 1$; then $a, b \le 1$, and we use (C). Then $v$ has degree $2$ in $K$, and $v$ is not a terminal of $L$: if $a = b = 1$ because $u$ is not a terminal of $M$, and if $a = 0$ and $b = 1$ because otherwise $u$ would also be a terminal of $L$, and $\{u, v\}$ would be a block of $L$ and $\{u\}$ a block of $M$, contrary to $b = 1$. This proves (*).}\fixnote{fix:11.1-claim2-proof-added}{The paper takes for granted that $G/e$ is an intertwining. The argument above is the one used in the formalization.} From the inductive hypothesis, $G/e$ has rank $\le r(p)$; but it has the same rank as $G$ (for $e$ is not a loop since $e \in E(L_1)$) and so $G$ has rank $\le r(p)$ as required. This proves (2).

\hspace*{12pt}(3) \textit{We may assume that $L_1$ and $L_2$ are vital, and $E(L_1 \cap L_2) = \emptyset$.}

\textit{Subproof.} Suppose that there is a linkage $L'_1$ in $G$ with the same pattern as $L_1$, with $E(L'_1 \cap L_2) \ne \emptyset$. Let $e \in E(L'_1 \cap L_2)$; then $G/e$ is an intertwining of the $p$-linkages $L'_1/e$ and $L_2/e$, and the result follows from the inductive hypothesis. We may therefore assume that there is no such $L'_1$. Since $L'_1 \cup L_2 = G$ for every linkage $L'_1$ with the same pattern as $L'_1$, it follows that $E(L'_1) = E(G) \setminus E(L_2)$, and in particular $E(L_1) = E(G) \setminus E(L_2)$. Since $E(L'_1) = E(L_1)$ and hence $L'_1 = L_1$ for every linkage $L'_1$ with the same pattern as $L_1$, it follows from (2) that $L_1$ is vital, and similarly we may assume that $L_2$ is vital. This proves (3).

From (1.1) and (3), $G$ has tree-width $\le w$. From (11.2) there is a linked tree-decomposition $(T, W)$ of $G$ of width $\le w$, such that every vertex of $T$ has degree $\le 3$\fix{, and we may assume that no edge of $G$ belongs to $W(t)$ for two distinct $t \in V(T)$, by deleting edges from the graphs $W(t)$ as in the proof of (11.2)}.\fixnote{fix:11.1-proof}{This property is needed in Claim (5) below, where $(A_i, B_i) = (W(S_i), W(T_i))$ must be a separation of $G$ (see the note in Section 11, page 43). Deleting the edges keeps the sets $W(e)$, so $(T, W)$ stays linked.} \fix{Choose such a $(T, W)$ with $|V(T)|$ as small as possible.} From (11.3)\fix{, since the tree $T'$ of (11.3) has fewer vertices than $T$ and $(T', W')$ has all the properties above,} we deduce\fixnote{fix:11.1-claim4}{As printed, Claim (4) does not follow. When (4) fails, (11.3) gives a smaller linked tree-decomposition, and this is a contradiction only if $|V(T)|$ was chosen as small as possible. For example, let $G$ have one vertex, let $T$ be a path $t_0 t_1 t_2$, and let every $W(t)$ be $G$. With $e_1 = t_0 t_1$ and $e_2 = t_1 t_2$, (4) fails. The tree $T'$ of (11.3) keeps maximum degree $\le 3$, since its new edge replaces $e_1$ and $e_2$ at their ends.}
